% !TeX root = se100_the_ontological_neutrality_theorem.tex
\section{Implications for Ontology Design}
\label{sec:implications}

The result established in Section~\ref{sec:formal-dev}
does not prescribe a particular ontology.
Rather, it constrains the design space of any ontology
intended to function as a neutral substrate
across admissible interpretive frameworks under persistent disagreement.
This section clarifies what such constraints require,
what they exclude, and
how interpretation, evaluation, and explanation may be accommodated
without violating neutrality.

The governing distinction is the one established in the previous sections:
a neutral substrate may represent that a claim was made,
by whom, and with what provenance,
but it may not assert the truth of
causal or normative claims as substrate-layer facts.
The exclusion is therefore at the level of assertion,
not vocabulary.
Causal and normative content may appear as reified,
attributed, provenance-bearing discourse;
it may not be asserted of substrate individuals
as a foundational ontological commitment.

%-----------------------------------------------------
\subsection{Substrate Exclusions}
\label{sec:substrate-exclusions}
%-----------------------------------------------------

The primary implication is negative but precise.
A neutral ontological substrate must exclude,
at the foundational layer,
assertions that are framework-variant across admissible frameworks.
In particular, the substrate
must not assert causal or normative conclusions as substrate-layer facts.

Normative exclusions include substrate-layer assertions
that an action, event, actor, institution, or artifact
was permitted, prohibited, obligatory, compliant, noncompliant,
justified, violative, blameworthy, liable, or responsible
in a framework-dependent sense.
Such conclusions presuppose legal, political, institutional,
or evaluative standards that may vary across admissible frameworks.

Causal exclusions include substrate-layer assertions
that one event caused, produced, prevented, enabled,
was responsible for, or explains another.
Such conclusions presuppose causal models,
counterfactual assumptions, variable selections,
mechanistic accounts, evidentiary thresholds,
or legal standards of causation that may vary
across admissible frameworks.

These exclusions do not deny the importance of causal or normative reasoning for accountability.
Rather, they preserve the distinction between representation
and interpretation.
Embedding causal or normative conclusions as substrate-layer facts
collapses that distinction,
commits the substrate to a particular interpretive framework,
and thereby violates neutrality under the theorem.

The exclusions apply specifically to asserted conclusions,
not to the referents about which such conclusions are made.
Actors, occurrences, institutional artifacts, jurisdictions,
records, documents, claims, reports, and provenance metadata
may all be represented without committing to how they are
causally explained or normatively evaluated.
The substrate may therefore be rich in referential structure
while remaining silent on contested interpretation.

%-----------------------------------------------------
\subsection{Substrate Requirements}
\label{sec:substrate-requirements}
%-----------------------------------------------------

Although the constraints are restrictive,
they do not leave the ontology empty or impoverished.
On the contrary, a neutral substrate must make strong,
explicit commitments in areas that
are invariant
across the admissible frameworks under consideration.
These commitments provide the shared referential ground
on which interpretive disagreement can occur.

In terms of the Method of Levels of Abstraction
(LoA)~\citep{floridi2008levels,floridi2011philosophy},
the substrate fixes a framework-invariant referential structure:
the predicate symbols and assertion forms needed to establish
referential regimes for entities, occurrences,
and institutional artifacts.
A referential regime comprises the substrate-level conditions
of individuation, co-reference, and persistence.
It determines what counts as one referent,
when two references denote the same referent,
and how that referent is tracked across time,
records, jurisdictions, or interpretive contexts.

The claim that a neutral substrate relies on such structure
does not assert that all questions of
individuation, co-reference, persistence, or boundary demarcation
are framework-independent.
Many such questions may be legitimately contested:
the spatial or temporal granularity of an occurrence,
the criteria for organizational continuity,
the recognition of informal institutions,
or the question whether two records refer to the same actor.
Where such referential questions are contested across
admissible frameworks,
the domain does not yet admit a neutral substrate
with respect to those frameworks.
The referential contestation must first be resolved or scoped out.

Neutrality instead requires that
the substrate fix a minimal and domain-appropriate referential regime
sufficient to support accountability relations.
These substrate-layer commitments are invariant
not because they resolve all ontological disagreement,
but because they remain stable under extension
by mutually incompatible causal, legal, analytic, or normative frameworks.

In particular, the substrate must provide:
\begin{itemize}
  \item stable reference to entities, occurrences,
        and institutional artifacts that participate
        in accountability relationships;
  \item explicit referential regimes specifying
        individuation, co-reference, and persistence conditions
        for those referents;
  \item provenance and attribution structures sufficient
        to record claims, reports, sources, documents,
        and interpretive contexts;
  \item structural relations such as temporal ordering,
        participation, containment, jurisdictional association,
        document linkage, and record membership,
        insofar as these assert neither causal efficacy
        nor normative force.
\end{itemize}

These commitments are not optional.
Accountability depends on the ability
to refer stably to who acted, what occurred,
which institutional instruments existed,
which records were made,
and within what jurisdictional or organizational scope.
When fixed as part of the shared substrate scope,
such commitments do not vary with interpretive stance
in the way causal or normative conclusions do,
and therefore do not threaten neutrality within that scoped domain.

%-----------------------------------------------------
\subsection{Externalizing Interpretation Without Loss}
\label{sec:externalizing-interpretation}
%-----------------------------------------------------

A common concern is that excluding causal and normative commitments
from the substrate renders the ontology incapable of supporting
explanation, evaluation, judgment, or accountability.
The result shows that the opposite is true:
such functions must be externalized in order to be supported robustly
across admissible disagreement.

Interpretation is accommodated by organizing the system
into distinct Levels of Organization (LoO)~\citep{floridi2011philosophy}.
The foundational LoO, the substrate,
records the referential ground:
entities, occurrences, institutional artifacts,
referential regimes, records, claims, attribution,
and provenance.
Higher-order LoOs record the interpretive layer:
causal explanations, legal classifications,
normative evaluations, responsibility assignments,
compliance judgments, or analytic conclusions.

By maintaining this structural separation,
multiple, mutually inconsistent interpretive layers
can operate over the same foundational substrate
without introducing inconsistency at the substrate layer.
Disagreement, revision, and contestation are represented
by the coexistence of multiple attributed interpretive assertions,
not by revision of substrate-layer facts.

Formally, the relation between the foundational layer
($\mathrm{LoO}_0$) and an interpretive layer
($\mathrm{LoO}_n$, where $n > 0$)
is a vertical stratification
governed by a read-only interface constraint.
An interpretive framework $\mathcal{F}$
layered at $\text{LoO}_n$
may expand the signature by defining
new causal, normative, legal, evaluative, or analytic predicates
over the referents established by $\mathrm{LoO}_0$.
However, this expansion constitutes a
conservative extension of the substrate:
the operational primitives of $\mathrm{LoO}_n$ are
strictly prohibited from mutating, retracting, or contradicting
the foundational constants, axioms, or identity mappings
asserted at $\mathrm{LoO}_0$.

This directional dependence guarantees that
changes or conflicts within higher-level interpretive logic
do not propagate downward to destabilize the substrate.
The substrate remains a stable point of reference;
the interpretive frameworks supply the competing explanations
and evaluations.

This architecture does not diminish the importance of causal reasoning
or normative judgment.
It relocates them to interpretive layers,
where the relevant assumptions are explicit,
attributed, contestable, and revisable.
The substrate does not certify compliance,
assign responsibility, or determine causation.
It makes such determinations adjudicable
by preserving the referents, claims, sources, and provenance
over which admissible frameworks may reason.

%-----------------------------------------------------
\subsection{Permitted Structural Relations and Excluded Commitments}
\label{sec:structural-relations}
%-----------------------------------------------------

The exclusion of causal and normative commitments from the substrate layer
does not entail the exclusion of all structured relations.
A neutral substrate permits relations that establish
referential, structural, temporal, documentary,
institutional, or provenance linkage
without asserting causal efficacy or evaluative force.

Examples of permitted substrate-layer relations may include:
\begin{itemize}
  \item temporal ordering, such as one occurrence preceding another;
  \item participation, such as an actor participating in an occurrence;
  \item containment or membership, such as a record belonging to a collection;
  \item provenance, such as a claim being recorded by a source;
  \item attribution, such as a framework, agent, or document making a claim;
  \item jurisdictional or institutional association,
        insofar as the relation identifies context
        without asserting legal consequence;
  \item co-reference relations, insofar as they belong
        to the fixed referential regime of the substrate.
\end{itemize}

By contrast, relations that embed causal attribution
or normative judgment are excluded from the substrate
as substrate-layer facts.
Examples include:
\begin{itemize}
  \item $\mathsf{Caused}(e_1,e_2)$ as a substrate assertion;
  \item $\mathsf{Responsible}(a,e)$ as a substrate assertion;
  \item $\mathsf{Violated}(a,r)$ as a substrate assertion;
  \item $\mathsf{Compliant}(s,r)$ as a substrate assertion;
  \item $\mathsf{Obligated}(a,\alpha)$ as a substrate assertion.
\end{itemize}

Such predicates may still appear as the opaque content
of reified assertions, reports, or framework-level conclusions.
For example, the substrate may assert:
[
$\mathsf{Asserts}(\mathcal{F}, \mathsf{Caused}(e_1,e_2))$.
]
This asserts that framework $\mathcal{F}$ made the causal claim.
It does not assert
$\mathsf{Caused}(e_1,e_2)$
as a substrate-layer fact.

The neutrality boundary is therefore drawn not at the presence
of structure,
nor at the syntactic occurrence of causal or normative vocabulary,
but at the embedding of framework-dependent conclusions
as foundational ontological facts.

%-----------------------------------------------------
\subsection{Reification as a Boundary Mechanism}
\label{sec:reification}
%-----------------------------------------------------

Reification plays a critical role in
enforcing the boundary between substrate and interpretation.
By treating claims, reports, descriptions, and assertions
as entities in their own right,
the ontology can represent discourse about the world
without asserting conclusions about the world itself.

However, reification functions as a boundary mechanism
only when applied consistently.
If causal or normative relations are asserted as
substrate-layer facts and merely duplicated as reified claims,
neutrality is not preserved.
The substrate has already endorsed the content.
Reification supports neutrality only when
the substrate refrains entirely from asserting the reified content as true.

Properly implemented, reification transforms
a framework-dependent proposition,
for example $\mathsf{Caused}(e_1,e_2)$,
into an assertion individual,
for example $A_{72}$.
The substrate may then assert framework-invariant metadata
about $A_{72}$:
\[
\mathsf{Assertion}(A_{72}), \quad
  \mathsf{RecordedBy}(A_{72}, \mathcal{F}_1), \quad
  \mathsf{HasContent}(A_{72}, \ulcorner \mathsf{Caused}(e_1,e_2) \urcorner).
\]
The quotation notation indicates that the causal proposition
is represented as content,
not asserted as a substrate-layer fact.

This type shift changes the representational status of the claim.
Instead of evaluating
$\mathsf{Caused}(e_1,e_2)$
directly against the world at the substrate layer,
the substrate evaluates only framework-invariant metadata
about an assertion object:
who made it, when it was recorded,
what it concerns,
what provenance attaches to it,
and under which interpretive context it was made.
The truth conditions of the claim remain inside the interpretive layer.

This mechanism preserves category discipline.
It distinguishes:
\begin{itemize}
  \item substrate-layer facts about assertion objects;
  \item framework-level conclusions contained within those assertions;
  \item provenance or attribution metadata about the source of the claim;
  \item later evaluations of the claim by admissible frameworks.
\end{itemize}

Thus, reification and contextual modeling do not evade
the theorem's constraint.
They satisfy it when they externalize interpretive commitments
rather than embedding them under a different description.
Neutrality is not achieved by representational cleverness,
but by disciplined exclusion of causal and normative commitments from
substrate-layer assertion.

%-----------------------------------------------------
\subsection{Accountability Preservation}
\label{subsec:accountability}
%-----------------------------------------------------

Externalizing causal and normative interpretation
does not weaken accountability.
It preserves accountability under disagreement
by allowing competing explanations and evaluations
to coexist over a stable shared referential base.

In accountability contexts,
different parties may reasonably disagree
about what caused an outcome,
whether a rule was violated,
whether an actor was responsible,
or whether an institutional response was justified.
If the substrate embeds one of these conclusions as fact,
then the record is no longer neutral with respect to the parties
who reject that conclusion.
The substrate becomes a site of adjudication
rather than a shared basis for adjudication.

A neutral substrate instead records
the referents, claims, sources, and provenance
needed for later evaluation.
It allows one framework to assert causation,
another to deny responsibility,
another to classify conduct as compliant,
and another to challenge that classification,
all while preserving stable reference
to the same entities, occurrences, records, and institutional artifacts.

The substrate therefore does not function as a truth oracle.
It does not decide whether a recorded claim is correct.
It guarantees something narrower and more durable:
that claims, sources, provenance, and referents
are faithfully recorded and stably available
for evaluation by admissible frameworks.
Truth-evaluation is externalized,
not abandoned.

This is the practical significance of the theorem.
A substrate that refuses to adjudicate causal or normative truth
at the foundational layer is not truth-indifferent.
It is designed so that truth, responsibility,
compliance, and causation can be evaluated
without the record itself having prejudged the matter.

%-----------------------------------------------------
\subsection{Summary of Structural Constraints}
\label{sec:summary-constraints}
%-----------------------------------------------------

The structural constraints established by the theorem can be summarized succinctly.
Any ontology satisfying the neutrality requirements,
under the contestability assumption and the referential-invariance
scope condition, must:

\begin{itemize}
  \item restrict the foundational LoA to assertions of
        framework-invariant referential structure,
        including the individuation, co-reference, and persistence
        conditions by which entities, occurrences, and institutional
        artifacts are fixed and tracked;
  \item assert no framework-variant proposition at the substrate layer,
        including no proposition whose truth conditions
        depend on causal or normative interpretation,
        regardless of current consensus among admissible frameworks;
  \item preserve the LoO hierarchy so that interpretation,
        evaluation, and explanation are supported exclusively
        through externalized higher-level organizational layers
        that extend but do not revise the substrate;
  \item satisfy extension stability by remaining consistent
        under extension by mutually incompatible interpretive frameworks,
        without revision of substrate-layer commitments.
\end{itemize}

These constraints do not define a particular ontology.
They delimit the structural form any ontology must take
if it is to function as a neutral substrate
for accountability under persistent interpretive disagreement.

The following section concludes by
situating these constraints as boundary conditions
on ontologies intended to function as
neutral substrates across domains and applications
in which the neutrality assumptions hold.
